\hypertarget{target:chap2}{}\chapter{مرور ادبیات}\label{chap:literature}
در سال‌های اخیر، مفهوم انتگرال و کاربردهای آن در شاخه‌های مختلف علمی مورد توجه بسیاری از پژوهشگران قرار گرفته است. انتگرال به‌عنوان یکی از ابزارهای اصلی در تحلیل مسائل پیوسته، نقش مهمی در حل مسائل مربوط به فیزیک، مهندسی و علوم کاربردی ایفا می‌کند.

\section{تاریخچه و مفاهیم اولیه انتگرال}
مفهوم انتگرال برای نخستین بار در قرن هفدهم توسط ریاضی‌دانانی مانند نیوتن\footnote{\lr{Newton}} و لایب‌نیتس\footnote{\lr{Leibniz}} معرفی شد. با گذشت زمان، تعریف دقیق‌تری از انتگرال ارائه شد. تعریف ریمان\footnote{\lr{Riemann}} از انتگرال، یکی از مهم‌ترین پیشرفت‌ها در این زمینه بود که امکان محاسبه مساحت به‌صورت حد مجموع‌های جزئی را فراهم کرد. اگر تابع $f(x)$ بر روی بازه $[a,b]$ تعریف شده باشد، انتگرال آن به‌صورت حد مجموع‌های ریمان در رابطه \eqref{eq:riemann} بیان می‌شود:
\begin{equation}
	\int_{a}^{b} f(x) \, dx = \lim_{n \to \infty} \sum_{i=1}^{n} f(x_i) \Delta x.
	\label{eq:riemann}
\end{equation}

\section{کاربردهای انتگرال در مطالعات پیشین}
در تحقیقات انجام‌شده، کاربردهای متنوعی برای انتگرال گزارش شده است. به‌عنوان مثال، در بسیاری از مطالعات مهندسی، از انتگرال برای محاسبه مساحت و حجم اجسام پیچیده استفاده شده است~\cite{first} همچنین، در حوزه فیزیک، انتگرال نقش مهمی در محاسبه کمیت‌هایی مانند کار، انرژی و چگالی ایفا می‌کند~\cite{second}

علاوه بر این، برخی پژوهش‌ها به بررسی روش‌های عددی برای محاسبه انتگرال پرداخته‌اند. از جمله این روش‌ها می‌توان به روش ذوزنقه‌ای\footnote{\lr{Trapezoidal Rule}} و سیمپسون\footnote{\lr{Simpson's Rule}} اشاره کرد که دقت مناسبی در تقریب دارند~\cite{third}